\documentclass[a4paper,12pt]{article}

\usepackage[T2A]{fontenc}
\usepackage[utf8]{inputenc}
\usepackage[russian]{babel}
\usepackage{geometry}
\usepackage{setspace}
\usepackage{hyperref}
\usepackage{enumitem}

\geometry{left=25mm,right=15mm,top=20mm,bottom=20mm}
\onehalfspacing
\setlist{nosep}

\begin{document}

\begin{titlepage}
    \begin{center}
        \textbf{МИНИСТЕРСТВО НАУКИ И ВЫСШЕГО ОБРАЗОВАНИЯ\\
        РОССИЙСКОЙ ФЕДЕРАЦИИ}\\[0.5cm]

        \textbf{Федеральное государственное автономное образовательное учреждение\\
        высшего образования\\
        «Национальный исследовательский университет ИТМО»}\\[1.5cm]

        Факультет программной инженерии и компьютерной техники\\
        Дисциплина: \textit{Веб-технологии}\\[3cm]

        {\Large \textbf{ОТЧЕТ ПО ЛАБОРАТОРНОЙ РАБОТЕ №3}}\\[0.4cm]
        {\large \textbf{«Разработка одностраничного веб-приложения (SPA)\\
        с использованием фреймворка Vue.js»}}\\[3cm]

        \begin{flushright}
            Выполнил: студент группы J3210\\
            Тищенко Павел\\[0.3cm]
            Проверил: \underline{\hspace{4cm}}
        \end{flushright}

        \vfill
        Санкт-Петербург\\
        2026
    \end{center}
\end{titlepage}

\tableofcontents
\newpage

\section{Цель работы}
Цель лабораторной работы --- миграция ранее разработанного приложения (ЛР1/ЛР2) в формат одностраничного веб-приложения с использованием Vue.js. В рамках выполнения необходимо показать:
\begin{itemize}
    \item подключение и использование клиентской маршрутизации;
    \item работу с API через библиотеку axios;
    \item разумное компонентное разбиение интерфейса;
    \item использование composable для повторно используемой логики.
\end{itemize}

\section{Постановка задачи}
Исходное приложение по тематике онлайн-курсов было реализовано ранее в формате многостраничного сайта. В ЛР3 требовалось:
\begin{enumerate}
    \item Перевести проект на архитектуру SPA.
    \item Сохранить пользовательские сценарии: каталог, карточка курса, авторизация, роли, обучение, личные кабинеты.
    \item Организовать отдельный mock API для ЛР3.
    \item Сохранить базовый и понятный интерфейс без избыточной визуальной сложности.
\end{enumerate}

\section{Теоретическая основа}
\subsection{Что такое SPA}
SPA (Single Page Application) --- это подход, при котором приложение загружается один раз, а дальнейшая навигация происходит внутри клиента за счет маршрутизатора и динамической подмены компонентов. В отличие от MPA, не требуется отдельная HTML-страница под каждый экран.

\subsection{Роль Vue.js в реализации SPA}
Vue.js предоставляет:
\begin{itemize}
    \item реактивную модель данных;
    \item компонентный подход;
    \item декларативную разметку;
    \item удобную интеграцию с экосистемой (router, store, HTTP-клиент).
\end{itemize}
Это позволяет последовательно наращивать функциональность без резкого усложнения проекта.

\subsection{Маршрутизация на клиенте}
Клиентский роутер обеспечивает отображение разных экранов в рамках одного приложения. Дополнительно используются guard-механизмы:
\begin{itemize}
    \item ограничение доступа к защищенным маршрутам без авторизации;
    \item разграничение доступа по ролям;
    \item запрет открытия страницы авторизации для уже вошедшего пользователя.
\end{itemize}

\subsection{Управление состоянием}
Для согласованности данных между страницами применяется централизованное состояние:
\begin{itemize}
    \item хранение данных пользователя и сессии;
    \item хранение коллекций курсов;
    \item минимизация дублирования логики загрузки и обновления данных.
\end{itemize}

\subsection{Работа с API}
Запросы вынесены в отдельный слой API. Такой подход позволяет:
\begin{itemize}
    \item разделить UI и сетевую логику;
    \item централизованно настраивать baseURL и заголовки;
    \item переиспользовать методы между страницами.
\end{itemize}

\subsection{Composable как механизм переиспользования}
Composable-файлы используются для повторяющейся функциональности (сессия, тема и т.д.). Это упрощает поддержку и уменьшает объем однотипного кода во view-компонентах.

\section{Архитектурное решение}
\subsection{Структура проекта}
Проект разделен на две части:
\begin{itemize}
    \item \textbf{frontend (spa)} --- Vue-приложение;
    \item \textbf{backend (server)} --- mock API с данными и базовыми правилами доступа.
\end{itemize}

\subsection{Frontend-уровни}
В клиентской части выделены:
\begin{itemize}
    \item \textbf{views} --- экраны приложения;
    \item \textbf{components/layouts} --- переиспользуемые элементы и общий каркас;
    \item \textbf{stores} --- состояние и бизнес-операции;
    \item \textbf{api} --- сетевой слой;
    \item \textbf{composables} --- общая логика, не привязанная к одному экрану.
\end{itemize}

\subsection{Backend-уровень}
Mock API содержит:
\begin{itemize}
    \item сущности пользователей и ролей;
    \item курсы, лекции, материалы, задания;
    \item записи обучения (enrollments), сертификаты, обсуждения;
    \item базовые ограничения доступа на операции изменения данных.
\end{itemize}

\section{Реализованный функционал}
\subsection{Базовые сценарии пользователя}
\begin{itemize}
    \item просмотр каталога с фильтрацией;
    \item просмотр страницы курса;
    \item регистрация и вход с учетом роли;
    \item покупка курса и переход к обучению;
    \item фиксация прогресса и выдача сертификата при завершении.
\end{itemize}

\subsection{Ролевые сценарии}
\begin{itemize}
    \item \textbf{Студент}: личный кабинет, курсы, прогресс, сертификаты;
    \item \textbf{Тренер}: личный кабинет, добавление курсов, лекций и семинаров.
\end{itemize}

\subsection{Навигация и защита маршрутов}
Навигация реализована через единый router. Для защищенных экранов применяется проверка авторизации и роли. Это предотвращает доступ к несоответствующим разделам напрямую по URL.

\subsection{Согласованность интерфейса}
Интерфейс намеренно оставлен базовым: минимальный CSS, акцент на читаемость и функциональность. Такое решение соответствует формату учебной работы и снижает риск регрессий при миграции.

\section{Проверка работоспособности}
Проведена ручная проверка основных сценариев:
\begin{itemize}
    \item запуск API и frontend в отдельных процессах;
    \item успешная регистрация и вход;
    \item корректная подгрузка каталога и карточек курсов;
    \item доступ к защищенным разделам только после авторизации;
    \item корректное разделение студент/тренер;
    \item обновление прогресса в обучении и генерация сертификата.
\end{itemize}

\section{Анализ результатов}
В ходе миграции подтверждено, что выбранная архитектура позволяет перенести прикладную логику из ЛР1/ЛР2 в SPA без потери ключевых сценариев. Разделение на маршруты, компоненты, store, composables и API-модули упростило дальнейшее развитие проекта.

С практической точки зрения получен рабочий каркас для последующего масштабирования: можно добавлять новые сущности и сценарии, не ломая базовый поток пользователя.

\section{Заключение}
Цель лабораторной работы достигнута: реализовано одностраничное приложение на Vue.js с клиентской маршрутизацией, API-взаимодействием через axios, компонентной структурой и использованием composable.

Полученный результат соответствует ключевым требованиям ЛР3 и демонстрирует переход от многостраничного подхода к современной SPA-архитектуре.

\end{document}
